% Generated draft: requires mathematical and source review.
\documentclass[11pt]{article}
\usepackage[margin=25mm]{geometry}
\usepackage{fontspec}
\setmainfont{Latin Modern Roman}
\usepackage{amsmath,amssymb,mathtools,unicode-math}
\setmathfont{Latin Modern Math}
\usepackage{enumitem,needspace,xurl,hyperref}
\setlength{\parindent}{0pt}
\setlength{\parskip}{5pt}
\setlength{\emergencystretch}{3em}
\newcommand{\N}{\mathbb{N}}
\newcommand{\Z}{\mathbb{Z}}
\newcommand{\Q}{\mathbb{Q}}
\newcommand{\R}{\mathbb{R}}
\newcommand{\C}{\mathbb{C}}
\newcommand{\F}{\mathbb{F}}
\newcommand{\A}{\mathbb{A}}
\newcommand{\PP}{\mathbb{P}}
\newcommand{\E}{\mathbb{E}}
\newcommand{\eps}{\varepsilon}
\newcommand{\dd}{\,\mathrm{d}}
\begin{document}
\section*{UnsolvedMath ID: 801187}
\textbf{Source problem code:} catalogue entry
\subsection{Definitions and mathematical statement}
Let \(k\geq 1\) be an integer, and let \(\mu\subseteq\lambda\) be partitions. Their skew Young diagram is
\[
\lambda/\mu=\{(i,j):\mu_i<j\leq\lambda_i\}.
\]
A \(k\)-ribbon is an edge-connected skew Young diagram consisting of \(k\) cells and containing no \(2\times2\) block. Assume that \(\lambda/\mu\) admits a tiling by \(k\)-ribbons, and put \(n=(|\lambda|-|\mu|)/k\).

The supplied record refers to LLT \(k\)-ribbon tableaux but does not specify their labeling and semistandard admissibility conditions. Denote the intended family by \(\mathcal T^{(k)}_{\lambda/\mu}\). For a tableau \(T\) whose ribbons carry positive integer labels, let \(m_i(T)\) count the ribbons labeled \(i\), and set
\[
x^T=\prod_{i\geq1}x_i^{m_i(T)},
\qquad
H^{(k)}_{\lambda/\mu}(x;q)
=\sum_{T\in\mathcal T^{(k)}_{\lambda/\mu}}q^{s(T)}x^T.
\]
Here \(s(T)\) is the spin statistic intended by the source. Its normalization is not specified. For example, if the height of a ribbon is the number of rows it occupies, the sum of ribbon heights minus one and half of that sum give different exponent conventions. The latter can introduce half-integral powers of \(q\). The intended convention must therefore be fixed before interpreting polynomiality in \(q\).

\textbf{Requested statement, subject to these unresolved conventions.}
For every admissible \(k\) and shape, establish the following assertions for the intended LLT ribbon generating function:
\begin{enumerate}
\item[(a)] Give a combinatorial proof of symmetry: for every permutation \(\sigma\) of the positive integers that fixes all but finitely many indices,
\[
H^{(k)}_{\lambda/\mu}(x_{\sigma(1)},x_{\sigma(2)},\ldots;q)
=H^{(k)}_{\lambda/\mu}(x_1,x_2,\ldots;q).
\]
\item[(b)] Writing its homogeneous degree-\(n\) Schur expansion as
\[
H^{(k)}_{\lambda/\mu}(x;q)
=\sum_{\nu\vdash n}c^{(k)}_{\lambda/\mu,\nu}(q)s_\nu(x),
\]
where \(s_\nu\) is the Schur symmetric function, determine whether
\[
c^{(k)}_{\lambda/\mu,\nu}(q)\in\mathbb Z_{\geq0}[q]
\qquad\text{for every }\nu\vdash n.
\]
\end{enumerate}
The skew notation above makes the opening hypothesis explicit. The supplied statement subsequently uses \(H^{(k)}_\lambda\) and tableaux of shape \(\lambda\), so it does not settle whether the intended scope includes arbitrary skew shapes or only straight shapes \((\mu=\varnothing)\). This scope requires confirmation.

The registry labels the problem partially solved. Its supplied background reports that part (b) was proved using Kazhdan--Lusztig theory of affine Hecke algebras, attributing this to Leclerc--Thibon (2000) and Grojnowski--Haiman (2007), while direct combinatorial bijective proofs of symmetry and general Schur positivity remain partially open. This is a supplied literature assessment, not an independently verified conclusion; its exact scope and conventions require review.
\subsection{Short English statement}
For LLT generating functions of labeled ribbon tableaux, seek a combinatorial proof that permuting the variables leaves the function unchanged, and ask whether every Schur coefficient is a polynomial in \(q\) with nonnegative integer coefficients. The supplied record reports a proof of the positivity assertion but retains combinatorial proof questions. Tableau admissibility, spin normalization, and the intended straight or skew shape scope remain unspecified.
\subsection{Sources}
\begin{itemize}
\item Supplied local UnsolvedMath registry record: permanent ID \(801187\), exact source problem code \texttt{catalogue entry}, in the AMR Open Problem Lists collection. The record attributes the problem to Stanley (1999) and labels it \texttt{partially\_solved}.
\item Stanley, \emph{Positivity Problems and Conjectures in Algebraic Combinatorics} (1999), Problem 14. Supplied source URL: \url{https://math.mit.edu/~rstan/papers/problems.pdf}. The supplied provenance records access on 2026-10-16 and extraction labeled \texttt{pdf+html\_verified}; these provenance assertions were not independently checked for this draft.
\item Supplied background attribution: Leclerc--Thibon (2000) and Grojnowski--Haiman (2007), for proofs of part (b) using Kazhdan--Lusztig theory of affine Hecke algebras. No titles, publication details, or links for these works are supplied.
\item The registry record was updated on 2026-10-16. Its accompanying literature assessment is supplied dated curation, not independently verified here. It reports that direct combinatorial bijective proofs of symmetry and general Schur positivity remain partially open, without specifying the remaining cases.
\end{itemize}
\end{document}
